\chapter{محیط‌های ریاضی و متنی}
\label{chap:envs}

در این قالب از پکیج پیشرفته \E{tcolorbox} برای طراحی محیط‌های زیبا استفاده شده است. سیستم ارجاع‌دهی هوشمند (\E{cleveref}) نیز پیاده‌سازی شده است که به سادگی می‌توانید به محیط‌ها ارجاع دهید (مانند \cref{def:ai} و \cref{thm:fundamental}).

\section{محیط‌های ریاضی (قضیه، لم، نتیجه)}
محیط‌های ریاضی با رنگ آبی متمایز شده‌اند تا از سایر متون جدا شوند.

\begin{theorem}[قضیه اساسی حساب دیفرانسیل و انتگرال]
\label{thm:fundamental}
فرض کنید تابع $f$ در بازه بسته $[a,b]$ پیوسته باشد. در این صورت تابع $F$ تعریف شده به صورت:
\[ F(x) = \int_a^x f(t) dt \]
در این بازه پیوسته و در بازه باز $(a,b)$ مشتق‌پذیر است و $F'(x) = f(x)$.
\end{theorem}

\begin{lemma}
\label{lem:sample}
اگر تابعی در نقطه‌ای مشتق‌پذیر باشد، در آن نقطه پیوسته نیز هست.
\end{lemma}

\section{محیط‌های توصیفی (تعریف، مثال، نکته)}
محیط‌های تعاریف و نکات با رنگ‌های سبز و زرد مشخص شده‌اند تا توجه خواننده را جلب کنند.

\begin{definition}[هوش مصنوعی]
\label{def:ai}
هوش مصنوعی \lr{(Artificial Intelligence)} شاخه‌ای از علوم رایانه است که هدف آن خلق ماشین‌هایی است که قادر به انجام وظایفی هستند که نیازمند هوش انسانی است.
\end{definition}

\begin{example}
\label{ex:sample}
این یک مثال برای نشان دادن نحوه کارکرد محیط‌های شماره‌دار است. همانطور که می‌بینید شماره‌گذاری‌ها بر اساس فصل انجام می‌شود.
\end{example}

\begin{note}
برای ارجاع دادن به بخش‌ها، شکل‌ها و قضیه‌ها، فقط کافیست از دستور \E{\textbackslash cref\{label\}} استفاده کنید.
\end{note}

\begin{warning}
\label{warn:sample}
این یک هشدار است! دقت کنید که ترتیب فراخوانی پکیج‌ها در فایل تنظیمات بسیار مهم است.
\end{warning}

\section{محیط‌های بدون شماره}
برخی از محیط‌ها مانند هشدار، حل و اثبات به صورت پیش‌فرض شماره‌گذاری نمی‌شوند.

\begin{question}
تابع $f(x) = x^3 + 2x$ داده شده است. مشتق این تابع را به دست آورید و مقدار آن را در $x=1$ محاسبه کنید.
\end{question}
\begin{solution}
مشتق تابع برابر با $f'(x) = 3x^2 + 2$ خواهد بود. بنابراین مقدار آن در $x=1$ برابر است با ۵.
\end{solution}

\begin{proof}[اثبات لم مشتق‌گیری]
این یک محیط اثبات است که برای نمایش مراحل استدلال‌های ریاضی طراحی شده و در انتهای آن یک مربع توپر قرار می‌گیرد.
\end{proof}
